% -------------------------------------------------
% VERSION TRACKER — No modificar este bloque
% Versión:    Tesis Humanizada (recortada A16)
% Archivo:    Final/Conclusiones.tex
% Última mod: 2026-09-19 (recorte 4697 → ~1500 palabras)
% Audiencia:   general
% -------------------------------------------------

\section{Conclusiones}

Pymetory nació de una pregunta concreta: ¿puede una pequeña empresa colombiana tener un sistema de inventarios que entienda preguntas en español y responda con información real de su bodega, sin depender de software costoso ni de conocimientos técnicos especializados? La respuesta, después de diseñar, implementar y probar el sistema a lo largo de 16 semanas de trabajo metódico, es afirmativa. Este capítulo sintetiza el cumplimiento de los objetivos, el estado de las reglas inmutables, los hallazgos técnicos y las lecciones metodológicas; el detalle completo del desarrollo multi-agente y de cada auditoría queda registrado en los anexos y en la documentación técnica del repositorio.

\subsection{Cumplimiento de Objetivos}

El proyecto alcanzó todos los objetivos que se propuso en la fase de formulación, cada uno verificado contra los indicadores de cumplimiento establecidos:

\begin{itemize}
    \item \textbf{Objetivo General}: Se diseñó e implementó un sistema web funcional para la gestión de inventario de materia prima, con un módulo de consulta en lenguaje natural (RAG) sobre datos reales de la base de datos. El sistema integra Laravel 13, React 19, Inertia.js 2.0 y MySQL 8.0, se despliega en la nube y cuenta con nueve módulos operativos documentados.

    \item \textbf{Objetivo Específico 1 -- Identificar necesidades (Cumplido)}: La entrevista semiestructurada con la administradora de una panificadora del Valle del Cauca permitió identificar el problema real de la empresa y derivar 18 requerimientos funcionales y 4 no funcionales (Capítulo \ref{chap:requerimientos}); de los funcionales, 13 quedaron implementados, 3 de forma parcial y 2 se proponen como trabajo futuro.

    \item \textbf{Objetivo Específico 2 -- Diseñar la arquitectura (Cumplido)}: Arquitectura MVC sobre Laravel 13 con Inertia.js 2.0 hacia React 19, base de datos de 10 tablas normalizadas que soporta FEFO, Kardex inmutable, RBAC y el módulo RAG, con diagramas de secuencia de los 4 procesos críticos. Indicador cumplido: documento de diseño completo aprobado.

    \item \textbf{Objetivo Específico 3 -- Implementar el módulo LLM (Cumplido)}: El \texttt{ChatLLMController} implementa el patrón RAG en tres métodos (\texttt{classifyQuery}, \texttt{buildRagContext}, \texttt{ask}), con clasificación de 10 categorías de intención, contexto real desde la base de datos y degradación elegante ante fallo del servicio de IA. Indicador cumplido: sistema funcional con 9 módulos y despliegue en la nube.

    \item \textbf{Objetivo Específico 4 -- Validar el sistema (Cumplido)}: 111 casos de prueba en 13 suites con una tasa de éxito del 99.1\%; 7 escenarios FEFO; precisión del clasificador RAG del 98\%; cumplimiento del nivel AA de WCAG 2.1. Indicador cumplido: 111 casos ejecutados con éxito superior al umbral del 95\%.
\end{itemize}

\subsection{Reglas Inmutables: Verificación}

Las cuatro condiciones impuestas por el director (Prof.~Héctor Fabio Ocampo) desde las primeras sesiones ---las ``reglas inmutables''--- quedaron verificadas de forma independiente:

\begin{enumerate}
    \item \textbf{FEFO obligatorio}: implementado como scope del modelo \texttt{Lote}, se aplica en cada despacho sin intervención manual. Las 7 pruebas FEFO confirmaron la priorización correcta del vencimiento más próximo incluso cuando un despacho consume varios lotes; el umbral de alerta temprana (15 días configurables) ataca la causa raíz identificada en la panificadora.
    \item \textbf{Kardex histórico inmutable}: triggers de base de datos que rechazan UPDATE/DELETE sobre \texttt{movimientos} y ausencia de rutas de modificación en la aplicación. Cada registro guarda usuario, marca de tiempo, cantidades antes/después y razón ---la trazabilidad que el administrador demandaba---.
    \item \textbf{Seguridad RBAC}: el middleware \texttt{CheckRole} centraliza el control de acceso; los roles Admin/Operario se probaron en 5 casos de autenticación y la interfaz oculta los módulos no autorizados en lugar de mostrar errores 403.
    \item \textbf{RAG Dinámico}: cada consulta adversa las tablas reales del inventario antes del prompt, eliminando el riesgo de alucinaciones sobre datos inexistentes; el mecanismo de fallback garantiza que, sin servicio de IA, el usuario siempre reciba los datos de su inventario.
\end{enumerate}

\subsection{Hallazgos Técnicos y de Metodología}

El desarrollo produjo contribuciones técnicas documentadas que trascienden el caso concreto:

\begin{itemize}
    \item \textbf{Inertia.js como alternativa a API REST}: eliminó endpoints separados, serialización manual y sincronización de tipos entre TypeScript y PHP. Relevantes para equipos pequeños y aplicaciones web de alcance moderado sin API pública.
    \item \textbf{Clasificación de intenciones con regex}: 98\% de precisión sobre 50 consultas en español colombiano con costo computacional nulo. La única falta (un lote mal clasificado) delimita con precisión el alcance del enfoque y motiva la búsqueda semántica como mejora futura.
    \item \textbf{Fallback como garantía de disponibilidad}: en los 4 escenarios de fallo simulados (sin API key, LLM inactivo, timeout, error de red) el sistema nunca quedó inoperativo. El principio ``la IA es valor agregado, no dependencia crítica'' es transferible a cualquier sistema que integre LLM en funciones operativas.
    \item \textbf{Infraestructura a costo cero es viable}: el sistema operó sobre los niveles gratuitos de OCI (instancia ARM de 24 GB) y GCP. La única limitación es la alta demanda de instancias ARM gratuitas de Oracle, mitigada con un script de sondeo automatizado.
    \item \textbf{Documentación concurrente}: registrar cada decisión al momento de tomarla convirtió la redacción de la tesis en compilación y refinamiento, no en reconstrucción.
\end{itemize}

\subsection{Impacto del Proyecto}

\begin{itemize}
    \item \textbf{Impacto tecnológico}: primer sistema de gestión de inventarios de código abierto desarrollado en Colombia que integra un módulo RAG de consulta en español. Demuestra que no se necesita entrenar un modelo propio para obtener respuestas precisas en un dominio cerrado y predecible.
    \item \textbf{Impacto económico y social}: elimina los costos de licencias propietarias (desde \$200.000 COP/mes por usuario); su costo anual de operación estimado (\$1.140.000 COP con hosting de gama baja, o solo el dominio si se despliega en OCI) es menor que un mes de Odoo para 5 usuarios. En un país donde solo el 42\% de las microempresas ha avanzado en transformación digital, ofrecer una alternativa gratuita, en español y con estándares del sector (FEFO) contribuye a cerrar la brecha de quienes hoy no tienen ninguna herramienta.
    \item \textbf{Impacto académico}: el proyecto deja un modelo de referencia documentado ---requisitos con PYME real, pruebas de usabilidad y accesibilidad, integración de IA--- para futuros trabajos de grado en ingeniería de sistemas.
\end{itemize}

\subsection{Lecciones Aprendidas}

\begin{itemize}
    \item \textbf{Desarrollo full-stack}: integrar Laravel, React, Inertia, Tailwind y MySQL enseñó a tomar decisiones de arquitectura con criterios de mantenibilidad, evaluando contexto (recursos unipersonales, cronograma académico, usuario no técnico) en lugar de aplicar recetas generales.
    \item \textbf{Integración de IA en aplicaciones reales}: el reto no es el modelo (que se consume como servicio), sino la ingeniería de software para integrarlo de forma robusta: clasificar intención, estructurar contexto, degradar con elegancia.
    \item \textbf{Pruebas automatizadas como red de seguridad}: cada sprint terminaba con la suite ejecutada; la inversión (~25\% del tiempo) se recuperó con creces en la fase de correcciones.
    \item \textbf{Desarrollo asistido por IA}: tres agentes en paralelo (OpenCode Go/DeepSeek, Claude Code/Sonnet, Gemini Antigravity) actuaron como amplificadores de capacidad; la experiencia delimitó qué tareas se benefician de la IA y cuáles exigen criterio humano. Se documenta como contribución metodológica.
\end{itemize}

\subsection{Logros del Desarrollo Multi-Agente}

La adopción de agentes de IA especializados bajo supervisión humana produjo resultados verificables: 6 agentes con dominios acotados (QA, backend, frontend, devops, documentación, repositorio); 35 + 42 hallazgos de auditoría frontend documentados con Playwright, 11 correcciones aplicadas que levaron Lighthouse accesibilidad de 62/100 a 91/100 y SEO de 58/100 a 89/100; más de 16 documentos técnicos generados (Delta); y una sincronización trazable entre casos de uso, issues de GitHub y agentes responsables. El detalle completo de cada ciclo de auditoría y corrección está en la documentación del repositorio y en el capítulo de pruebas.

De esa experiencia emergen cuatro lecciones metodológicas centrales (la bitácora completa está en los anexos):

\begin{enumerate}
    \item \textbf{La especialización supera a la generalidad}: agentes con dominios acotados (frontend, backend, QA) producen mejor calidad que un generalista, análogo a los equipos humanos.
    \item \textbf{La documentación concurrente habilita la coordinación}: el Kanban, el Plan Maestro y el contexto persistente funcionaron como memoria de largo plazo; sin ellos, los agentes tomaban decisiones inconsistentes entre sí.
    \item \textbf{La auditoría automatizada cubre puntos ciegos}:los aspectos no funcionales (rendimiento, accesibilidad, seguridad, SEO) que ningún desarrollador verifica manualmente en cada cambio --- el agente QA descubrió errores que ninguna revisión humana había encontrado.
    \item \textbf{Los agentes necesitan reglas de concurrencia explícitas}: asignar propiedad de archivos por agente y exigir criterios de finalización explícitos por ciclo eliminó los conflictos de escritura y los ciclos que se expandían sin control.
\end{enumerate}

\subsection{Conclusión Final}

Pymetory cumple con los objetivos que se propuso y con las reglas que se le impusieron. Los 111 casos de prueba (99.1\% de éxito), las 7 pruebas FEFO y las 50 evaluaciones del RAG confirman que el sistema funciona y entiende el español colombiano con datos reales.

Su valor trasciende las métricas al demostrar tres cosas: que la IA aplicada a problemas cotidianos del sector productivo no es un lujo inalcanzable; que la brecha digital de las PYMEs colombianas es de diseño para su contexto, no de falta de tecnología; y que un proyecto de grado puede ser el inicio de una herramienta que, con las mejoras trazadas en Trabajos Futuros, ayude a que las pequeñas empresas colombianas empiecen a tomar decisiones basadas en datos.
